%%%%%%%%%%%%%%%%%%%%%%%%%%%%%%%%%%%%%%%%%%%%%%%%%%%%%%%%%%%%%%%%%%%%%%%%%%%%%%%
%% Dados do trabalho
%%%%%%%%%%%%%%%%%%%%%%%%%%%%%%%%%%%%%%%%%%%%%%%%%%%%%%%%%%%%%%%%%%%%%%%%%%%%%%%

%% TÍTULO PROVISÓRIO (working title). Finalizar só após o ensaio do protótipo:
%% o título deve refletir o que foi realizado, e os resultados ainda podem mudar.
%%   - Descritivo, sem afirmar resultado ("limites" entra só depois do ensaio).
%%   - Sem nome de projeto (Kaelix) nem modelo comercial (MPU6050, ESP32-S3):
%%     a NBR 14724 exige que o título permita indexação e recuperação.
%%   - Sem sigla (MEMS): desaconselhada em título por IEEE e Nature.
%%   - O subtítulo descreve o conteúdo; na versão final, é onde entra a
%%     contribuição apurada.
%% A quebra depois dos dois-pontos separa título e subtítulo nas três folhas; nos
%% metadados do PDF ela vira espaço.
\titulo{Monitoramento de vibração de baixo custo para manutenção preditiva:\texorpdfstring{\\}{ }aquisição e detecção embarcada de anomalias em motores}
\autor{Lauan Matheus da Rocha Alves}
\local{Teresina -- PI}
\data{2026}
\orientador{Lia Mariana Nascimento Brito}
%\coorientador{Nome do coorientador}
\preambulo{Trabalho de Conclusão de Curso apresentado como requisito parcial
para obtenção do título de Bacharel em Engenharia de Software do Instituto de
Ensino Superior -- iCEV.}

%% Metadados do PDF: precisam vir depois de \titulo e \autor.
\makeatletter
\hypersetup{pdftitle = {\@title}, pdfauthor = {\@author}}
\makeatother

%% Usados na capa do modelo iCEV (preambulo/estilo-icev.tex).
\newcommand{\cursoicev}{Curso de Bacharelado em Engenharia de Software}
\newcommand{\instituicaoicev}{Instituto de Ensino Superior -- iCEV}
